\chapter{基本概念 (Basic Concepts)}
\label{chapter:basicConcepts}



%----------NOTATION
\section{关于符号的说明 (A few words about notation)}

本报告的一个核心目标是收集、审查、纠正和统一关于 Monero 加密货币内部运作的所有现有信息，同时提供所有必要的细节，以循序渐进的方式呈现材料。

实现这一目标的一个重要工具是确定一套符号约定。其中，我们使用了：

\begin{itemize}
\item 小写字母表示简单值、整数、字符串、位表示等
\item 大写字母表示曲线点和复杂结构
\end{itemize}

对于具有特殊含义的元素，我们尽量在全文中使用相同的符号。例如，曲线生成元始终用 \(G\) 表示，其阶为 \(l\)，私钥/公钥尽可能分别用 \(k/K\) 表示，等等。
\\

除此之外，我们在呈现算法和方案时力求{\em 概念化}。具有计算机科学背景的读者可能会觉得我们忽略了诸如元素的位表示等问题，或在某些情况下忽略了如何执行具体操作。此外，数学专业的学生可能会发现我们忽略了抽象代数的解释。

然而，我们并不认为这是一种损失。简单的对象如整数或字符串总是可以用位串表示。所谓"字节序"很少相关，对我们的算法来说主要是约定问题。\footnote{在计算机内存中，每个字节存储在自己的地址中（地址类似于一个编号的插槽，字节可以存储在其中）。一个给定的"字"或变量通过其字节的最低地址来引用。如果变量 $x$ 有4个字节，存储在地址10-13中，则使用地址10来查找 $x$。$x$ 的字节在其地址集中的组织方式取决于{\em 字节序}，尽管每个单独的字节在其地址内始终以相同方式存储。基本上，$x$ 的哪一端存储在引用地址中？它可以是{\em 大端}或{\em 小端}。给定 $x = $ 0x12345678（十六进制，2个十六进制数字占用1个字节即8个二进制位），以及地址数组 \\{10, 11, 12, 13\\}，$x$ 的大端编码为 \\{12, 34, 56, 78\\}，小端编码为 \\{78, 56, 34, 12\\}。\\cite{endianness}}

椭圆曲线点通常用数对 \((x, y)\) 表示，因此可以用两个整数来表示。然而，在密码学领域中，通常会应用{\em 点压缩}技术，允许仅使用一个坐标的空间来表示一个点。对于我们的概念化方法来说，是否使用点压缩通常是次要的，但在大多数情况下是隐式假定的。\\

我们\marginnote{src/crypto/ {\tt keccak.c}}还自由地使用了密码学 hash 函数，而未指定任何具体算法。在 Monero 的情况下，通常使用的是 \(\mathit{Keccak}\)\footnote{\label{kekkak_note}Keccak hash 算法构成了 NIST 标准 {\em SHA-3} \cite{nist-sha3} 的基础。}变体，但如果未明确提及，则对理论并不重要。

密码学 hash 函数（以下简称为"hash 函数"或"hash"）接受任意长度的消息 $\mathfrak{m}$，并返回固定长度的 hash $h$（或"消息摘要"），对于给定输入，每个可能输出的概率相等。密码学 hash 函数难以逆转，具有一个被称为{\em 大雪崩效应}的有趣特性，该特性可以使非常相似的消息产生非常不同的 hash，并且很难找到两个具有相同消息摘要的消息。

hash 函数将应用于整数、字符串、曲线点或这些对象的组合。这些情况应被解释为位表示的 hash，或此类表示的拼接。根据上下文，hash 的结果可以是数字、位串，甚至曲线点。这方面的更多细节将在需要时给出。



%----------MODULAR ARITHMETIC
\section{模运算 (Modular arithmetic)}
\label{sec:modular-arithmetic}

大多数现代密码学始于模运算，而模运算又始于取模运算（用 `mod' 表示）。我们只关心正取模，它始终返回一个正整数。

正取模类似于两个数相除后的"余数"，例如 $c$ 是 $a/b$ 的"余数"。让我们想象一条数轴。要计算 $c = a \pmod b$，我们站在点 $a$ 处，然后以每步 $\text{step} = b$ 的步伐向零走去，直到到达一个 $\geq{0}$ 且 $<b$ 的整数。那就是 $c$。例如，$4$（取模 3）$= 1$，$-5 \pmod 4 = 3$，依此类推。

形式化地，正取模在此定义为 $c = a \pmod b$ 满足 $a=bx+c$，其中 $0\leq{c}<{b}$，$x$ 是一个会被丢弃的带符号整数（$b$ 是正非零整数）。

注意，如果 $a \leq n$，则 $-a \pmod n$ 等于 $n - a$。


\subsection{模加法和模乘法 (Modular addition and multiplication)}
\label{subsec:modular-addition-multiplication}

在计算机科学中，进行模运算时避免大数是很重要的。例如，如果我们有 $29+87 \pmod{99}$，且不允许使用三位或更多位的变量（如 $116 = 29+87$），那么我们无法直接计算 $116 \pmod{99} = 17$。

要执行 $c = a+b \pmod n$，其中 $a$ 和 $b$ 各自小于模数 $n$，我们可以这样做：
\begin{itemize}
	\item 计算 $x = n-a$。如果 $x > b$ 则 $c = a+b$，否则 $c = b - x$。
\end{itemize}

我们可以利用模加法来实现模乘法（$a*b \pmod n = c$），使用的算法称为"双倍加法"。让我们通过示例来演示。假设我们要做 $7*8 \pmod 9 = 2$。这等同于
\[7*8 = 8+8+8+8+8+8+8 \pmod 9\]

现在将其分成两两一组。
\[(8+8) + (8+8) + (8+8) + 8\]

再分一次，两两一组。
\[[(8+8) + (8+8)] + (8+8) + 8\]

$+$ 点运算的总次数从6次减少到4次，因为我们只需要找到一次 $(8+8)$。\footnote{双倍加法的效果在大数时变得明显。例如，对于 $2^{15} * 2^{30}$，直接加法需要大约 $2^{15}$ 次 $+$ 运算，而双倍加法只需要15次！}\\

双倍加法是通过将第一个数（"被乘数" $a$）转换为二进制来实现的（在我们的例子中，7 $\rightarrow$ [0111]），然后遍历二进制数组进行双倍和加法。

让我们建立一个数组 $A = [0111]$ 并按3,2,1,0索引。\footnote{这被称为"LSB 0"编号，因为最低有效位的索引为0。本章其余部分将使用"LSB 0"。这里的关键是清晰性，而非精确的约定。} A[0] = 1 是 A 的第一个元素，也是最低有效位。我们设置一个结果变量初始为 $r = 0$，设置一个求和变量初始为 $s = 8$（更一般地，我们从 $s = b$ 开始）。我们遵循以下算法：
\begin{enumerate}
	\item 遍历：$i = (0,...,A_{size} - 1)$
	\begin{enumerate}
		\item 如果 A[i] == 1，则 $r = r + s \pmod n$。
		\item 计算 $s = s + s \pmod n$。
	\end{enumerate}
	\item 使用最终的 $r$：$c = r$。
\end{enumerate}

在我们的例子 $7*8 \pmod 9$ 中，出现如下序列：
\begin{enumerate}
	\item $i = 0$
	\begin{enumerate}
		\item A[0] = 1，所以 $r = 0 + 8 \pmod 9$ = 8
		\item $s = 8 + 8 \pmod 9$ = 7
	\end{enumerate}
	\item $i = 1$
	\begin{enumerate}
		\item A[1] = 1，所以 $r = 8 + 7 \pmod 9$ = 6
		\item $s = 7 + 7 \pmod 9$ = 5
	\end{enumerate}
	\item $i = 2$
	\begin{enumerate}
		\item A[2] = 1，所以 $r = 6 + 5 \pmod 9$ = 2
		\item $s = 5 + 5 \pmod 9$ = 1
	\end{enumerate}
	\item $i = 3$
	\begin{enumerate}
		\item A[3] = 0，所以 $r$ 保持不变
		\item $s = 1 + 1 \pmod 9$ = 2
	\end{enumerate}
	\item $r = 2$ 即为结果
\end{enumerate}


\subsection{模幂运算 (Modular exponentiation)}

显然 $8^7 \pmod 9 = 8*8*8*8*8*8*8 \pmod 9$。正如双倍加法一样，我们可以使用平方乘法。对于 $a^e \pmod{n}$：
\begin{enumerate}
	\item 定义 $e_{scalar} \rightarrow e_{binary}$；$A = [e_{binary}]$；$r = 1$；$m = a$
	\item 遍历：$i = (0,...,A_{size} - 1)$
	\begin{enumerate}
		\item 如果 A[i] == 1，则 $r = r * m \pmod n$。
		\item 计算 $m = m * m \pmod n$。
	\end{enumerate}
	\item 使用最终的 $r$ 作为结果。
\end{enumerate}


\subsection{模乘逆 (Modular multiplicative inverse)}

有时我们需要 $1/a \pmod n$，或者换句话说 $a^{-1} \pmod n$。某物乘以其逆元的定义是1（单位元）。想象 $0.25 = 1/4$，然后 $0.25*4 = 1$。\\

在模运算中，对于 $c = a^{-1} \pmod{n}$，$a c \equiv 1 \pmod{n}$ 且 $0 \leq c < n$，其中 $a$ 和 $n$ 互素。\footnote{在方程 $a \equiv b \pmod{n}$ 中，$a$ 与 $b$ 模 $n$ {\em 同余}，这仅意味着 \(a \pmod{n} = b \pmod{n}\)。}互素意味着它们除了1以外没有其他公因数（分数 $a/n$ 不能被约分/简化）。

当 $n$ 是素数时，我们可以使用平方乘法来计算模乘逆元，这基于{\em 费马小定理}：\footnote{\label{inverse_rule_note}模乘逆元有一条规则：\\
{\em 如果 $a c \equiv b \pmod{n}$ 且 $a$ 和 $n$ 互素，则该线性同余方程的解为 \(c = a^{-1} b \pmod{n}\)。}\cite{wiki-modular-arithmetic}\\
这意味着我们可以做 $c = a^{-1} b \pmod n \rightarrow ca \equiv b \pmod n \rightarrow a \equiv c^{-1} b \pmod n$。}\vspace{.175cm}
\begin{align*} 
    a^{n-1} &\equiv 1 \pmod{n} \\
    a*a^{n-2} &\equiv 1 \pmod{n} \\
    c \equiv a^{n-2} &\equiv a^{-1} \pmod{n}
\end{align*}

更一般地（也更快地），所谓的"扩展欧几里得算法" \cite{extended-euclidean} 也可以找到模逆元。


\subsection{模方程 (Modular equations)}
\label{subsec:modular-equations}

假设我们有一个方程 $c = 3*4*5 \pmod 9$。计算这个是直接的。给定两个表达式 $A$ 和 $B$ 之间的某种运算 $\circ$（例如 $\circ = *$）：\vspace{.175cm}
\[(A \circ B)\pmod{n} = {[A\pmod {n}] \circ [B\pmod{n}]}\pmod{n}\]

在我们的例子中，设 $A = 3*4$，$B = 5$，$n = 9$：\vspace{.175cm}
\begin{align*}
(3*4 * 5) \pmod{9} &= {[3*4 \pmod {9}] * [5 \pmod{9}]} \pmod{9} \\
				   &= [3]*[5] \pmod 9 \\
				c &= 6
\end{align*}

现在我们有了进行模减法的方法。\vspace{.175cm}
\begin{align*}
A - B \pmod n &\rightarrow A + (-B) \pmod n \\
			  &\rightarrow {[A \pmod {n}] + [-B \pmod{n}]} \pmod{n}
\end{align*}\\

同样的原理适用于像 $x = (a-b*c*d)^{-1} (e*f+g^{h}) \pmod n$ 这样的表达式。\footnote{大数的取模可以利用模方程。事实上 $254 \pmod {13} \equiv 2*10*10 + 5*10 + 4 \equiv (((2)*10 + 5)*10 + 4) \pmod {13}$。当 $a > n$ 时，$a \pmod n$ 的算法为：
\begin{enumerate}
	\item 定义 $A \rightarrow [a_{decimal}]$；$r = 0$
	\item 对于 $i = A_{size} - 1,...,0$
	\begin{enumerate}
		\item $r = (r*10 + A[i]) \pmod n$
	\end{enumerate}
	\item 使用最终的 $r$ 作为结果。
\end{enumerate}}



%----------ELLIPTIC CURVE CRYPTOGRAPHY
\section{椭圆曲线密码学 (Elliptic curve cryptography)}
\label{EllipticCurveCryptography}


\subsection{什么是椭圆曲线？ (What are elliptic curves?)}
\label{elliptic_curves_section}

一个\marginnote{{\tt fe}: 域元素}有限域 \(\mathbb{F}_q\)，其中 \(q\) 是大于3的素数，是由集合 \(\{0, 1, 2, ..., q-1\}\) 构成的域。加法和乘法 \((+,  \cdot)\) 以及取反 $(-)$ 是模 \(q\) 计算的。

``模 \(q\) 计算"意味着对两个域元素之间的任何算术运算实例，或对单个域元素的取反，都执行 \(\pmod q\)。例如，给定素域 \(\mathbb{F}_p\) 且 $p = 29$，$17+20=8$ 因为 $37 \pmod{29} = 8$。同样，$-13 = -13 \pmod{29} = 16$。\\

通常，具有给定 $(a,b)$ 数对的椭圆曲线定义为满足{\em Weierstraß}方程 \cite{Hankerson:2003:GEC:940321} 的所有坐标为 \((x, y)\) 的点的集合：\footnote{\label{notation1}符号说明：$a \in \mathbb{F}$ 表示 $a$ 是域 $\mathbb{F}$ 中的某个元素。}\vspace{.175cm}
\[y^2 = x^3 + a x + b \quad \textrm{where} \quad a, b, x, y \in \mathbb{F}_q\]

加密货币 Monero 使用一种属于所谓{\em 扭曲 Edwards}曲线 \cite{Bernstein2008} 类别的特殊曲线，通常表示为（对于给定的 $(a,d)$ 数对）：\vspace{.175cm}
\[a x^2 + y^2 = 1 + d x^2 y^2 \quad \textrm{where} \quad a, d, x, y \in \mathbb{F}_q \]


在下文中我们将偏好第二种形式。它相对于前面提到的 Weierstraß 形式的优势在于，基本密码学原语需要更少的算术运算，从而产生更快的密码学算法（详见 Bernstein {\em 等人} 在 \cite{Bernstein2007} 中的论述）。\\

设 \(P_1 = (x_1, y_1)\) 和 \(P_2 = (x_2, y_2)\) 是属于扭曲 Edwards 椭圆曲线（以下简称为 EC）的两个点。我们通过定义 $P_1 + P_2 = (x_1, y_1) + (x_2, y_2)$ 为点 $P_3 = (x_3, y_3)$ 来定义点加法，其中\footnote{通常椭圆曲线点在进行点加法等曲线运算之前会被转换为射影坐标\marginnote{src/crypto/ crypto\_ops\_ builder/ ref10Comm- entedComb- ined/\\\\ ge.h}，以避免执行域求逆运算以提高效率。\cite{ecc-projective}}\vspace{.175cm}
\begin{align*}
x_3 & =  \frac{x_1 y_2 + y_1 x_ 2}{1 + d x_1 x_2 y_1 y_2}  \pmod{q} \\
y_3 & =  \frac{y_1 y_2 - a x_1 x_2}{1 - d x_1 x_2 y_1 y_2} \pmod{q} 
\end{align*}

这些加法公式也适用于点倍乘；即当 \(P_1 = P_2\) 时。要减去一个点，将其坐标沿 y 轴取反，$(x,y) \rightarrow (-x,y)$ \cite{Bernstein2008}，然后使用点加法。回忆一下，$\mathbb{F}_q$ 中的"负"元素 $-x$ 实际上是 $-x \pmod{q}$。
\\

每当两个曲线点相加时，$P_3$ 是"原始"椭圆曲线上的一个点，换句话说，所有 $x_3,y_3 \in \mathbb{F}_q$ 并满足 EC 方程。

EC 中的每个点 $P$ 都可以使用自身的倍数从 EC 中的其他点生成一个阶（大小）为 $u$ 的子群。例如，某个点 $P$ 的子群可能有阶5，包含点 $(0, P, 2P, 3P, 4P)$，每个都在 EC 中。在 $5P$ 处，所谓的{\em 无穷远点}出现，它类似于 EC 上的"零"位置，坐标为 $(0, 1)$。\footnote{事实上，在上述加法运算下，椭圆曲线具有{\em 阿贝尔群}结构，因为它们的无穷远点是单位元。该概念的简洁定义可在 \url{https://brilliant.org/wiki/abelian-group/} 找到。}

方便的是，$5P + P = P$。这意味着该子群是{\em 循环的}。\footnote{\label{cyclical_note}循环子群意味着，对于 $P$ 的阶为 $u$ 的子群，以及任意整数 $n$，$n P = [n \pmod{u}] P$。我们可以想象自己站在地球仪上距某个"零"位置一定距离的一个点上，每走一步就移动那个距离。过一段时间后，我们会回到出发点，尽管可能需要绕很多圈才能精确回到原来的位置。回到完全相同位置所需的步数就是"步进群"的"阶"，我们所有的脚印都是该群中的唯一点。我们建议将此概念应用到这里讨论的其他想法中。}EC 中的所有 $P$ 都生成一个循环子群。如果 $P$ 生成的子群的阶是素数，则所有包含的点（除了无穷远点）都生成同一个子群。在我们的例子中，取点 $2P$ 的倍数：\vspace{.175cm}
\[2P, 4P, 6P, 8P, 10P \rightarrow 2P, 4P, 1P, 3P, 0\]

另一个例子：阶为6的子群 $(0, P, 2P, 3P, 4P, 5P)$。点 $2P$ 的倍数：\vspace{.175cm}
\[2P, 4P, 6P, 8P, 10P, 12P \rightarrow 2P, 4P, 0, 2P, 4P, 0\]
这里 $2P$ 的阶为3。由于6不是素数，并非所有成员点都能重建原始子群。

每个 EC 都有一个等于曲线上总点数（包括无穷远点）的阶 $N$，由点生成的所有子群的阶都是 $N$ 的因数（根据{\em 拉格朗日定理}）。我们可以想象所有 EC 点的集合 $\{0,P_1,...,P_{N-1}\}$。如果 $N$ 不是素数，某些点将生成阶等于 $N$ 因数的子群。

要找到任意给定点 $P$ 的子群的阶 $u$：
\begin{enumerate}
    \item 找到 $N$（例如使用 {\em Schoof 算法}）。
    \item 找到 $N$ 的所有因数。
    \item 对于 $N$ 的每个因数 $n$，计算 $n P$。
    \item 使得 $n P = 0$ 的最小 $n$ 即为子群的阶 $u$。
\end{enumerate} 

用于密码学的 EC\marginnote{{\tt ge}: 群元素} 通常具有 $N = hl$，其中 $l$ 是某个足够大的素数（如160位），$h$ 是所谓的{\em 余因子}，可以小到1或2。\footnote{具有小余因子的 EC 允许相对更快的点加法等操作。\cite{Bernstein2008}。}大小为 $l$ 的子群中的一个点通常被选为生成元 $G$，这是一种约定。对于该子群中的每个其他点 $P$，都存在一个满足 $P = n G$ 的整数 $0 < n \leq l$。

让我们扩展我们的理解。假设有一个阶为 $N$ 的点 $P'$，其中 $N=h l$。任何其他点 $P_i$ 都可以用某个整数 $n_i$ 找到，使得 $P_i=n_i P'$。如果 $P_1=n_1 P'$ 的阶为 $l$，那么任何阶为 $l$ 的 $P_2=n_2 P'$ 必须与 $P_1$ 在同一个子群中，因为 $l P_1=0 = l P_2$，并且如果 $l(n_1 P') \equiv l(n_2 P') \equiv N P'=0$，那么 $n_1$ 和 $n_2$ 必须都是 $h$ 的倍数。由于 $N= h l$，只有 $l$ 个 $h$ 的倍数，意味着只可能有一个大小为 $l$ 的子群。

简单来说，由 $(h P')$ 的倍数形成的子群始终包含 $P_1$ 和 $P_2$。此外，当 $n'$ 是 $l$ 的倍数时，$h(n' P')=0$，并且 $n' \pmod N$ 只有 $h$ 种这样的变化（包括 $n' = hl$ 时的无穷远点），因为当 $n' = h l$ 时它循环回到0：$h l P' = 0$。所以，EC 中只有 $h$ 个点 $P$ 使得 $h P$ 等于0。

类似的论证可以应用于任何大小为 $u$ 的子群。任何两个阶为 $u$ 的点 $P_1$ 和 $P_2$ 都在同一个子群中，该子群由 $(N/u) P'$ 的倍数组成。
\par
有了这个新的理解，显然我们可以使用以下算法来找到阶为 $l$ 的子群中的（非无穷远点）点：
\begin{enumerate}
    \item 找到椭圆曲线 EC 的 $N$，选择子群阶 $l$，计算 $h=N/l$。
    \item 在 EC 中选择一个随机点 $P'$。
    \item 计算 $P=h P'$。
    \item 如果 $P=0$ 则返回第2步，否则 $P$ 在阶为 $l$ 的子群中。
\end{enumerate}

计算\marginnote{{\tt sc}: 标量}任意整数 $n$ 和任意点 $P$ 之间的标量积 $nP$ 并不困难，而找到满足 $P_1 = n P_2$ 的 $n$ 则被认为是计算困难的。类似于模运算，这通常被称为{\em 离散对数问题}（DLP）。标量乘法可以被视为{\em 单向函数}，这为使用椭圆曲线进行密码学奠定了基础。\footnote{目前没有已知的方程或算法可以高效地（基于现有技术）求解 $P_1 = n P_2$ 中的 $n$，这意味着仅解开一个标量积就需要很多很多年。}

标量积 $nP$ 等价于 $(((P+P)+(P+P))…)$。虽然不总是最高效的方法，但我们可以像第 \ref{subsec:modular-addition-multiplication} 节中那样使用双倍加法。要获得总和 $R = n P$，请记住我们使用第 \ref{elliptic_curves_section} 节中讨论的 $+$ 点运算。

\begin{enumerate}
	\item 定义 $n_{scalar} \rightarrow n_{binary}$；$A = [n_{binary}]$；$R = 0$，即无穷远点；$S = P$
	\item 遍历：$i = (0,...,A_{size} - 1)$
	\begin{enumerate}
		\item 如果 A[i] == 1，则 R += S。
		\item 计算 S += S。
	\end{enumerate}
	\item 使用最终的 R 作为结果。
\end{enumerate}

注意，大小为 $l$ 的子群中的 EC 点（我们将从此使用）的标量是有限域 $\mathbb{F}_l$ 的成员。这意味着标量之间的算术运算取模 $l$。


\subsection{椭圆曲线公钥密码学 (Public key cryptography with elliptic curves)}
\label{ec:keys}
公钥密码学算法可以用类似于模运算的方式来设计。

设 \(k\) 为满足 \(0 < k < l\) 的随机选取的数，称之为{\em 私钥}。\footnote{私钥有时被称为{\em 密钥}。这使我们可以缩写：pk = 公钥，sk = 密钥。}通过标量积 \(k G = K\) 计算对应的{\em 公钥} \(K\)（一个 EC 点）。

由于{\em 离散对数问题}（DLP），我们无法轻易从 \(K\) 单独推导出 \(k\)。这一属性允许我们在标准公钥密码学算法中使用值 \((k, K)\)。


\subsection{椭圆曲线 Diffie-Hellman 密钥交换 (Diffie-Hellman key exchange with elliptic curves)}
\label{DH_exchange_section}

{\em Alice} 和 {\em Bob} 之间的基本 {\em Diffie-Hellman} \cite{Diffie-Hellman} 共享密钥交换可以按以下方式进行：

\begin{enumerate}
	\item Alice 和 Bob 各自生成自己的私钥/公钥对 \((k_A, K_A) \textrm{ 和 } (k_B, K_B)\)。双方公开或交换各自的公钥，并将私钥保留给自己。

	\item 显然，有 \[S = k_A K_B = k_A k_B G = k_B k_A G = k_B K_A\]

	Alice 可以私下计算 \(S = k_A K_B\)，Bob 计算 \(S = k_B K_A\)，从而使他们能够将这个单一值用作共享密钥。

	例如，如果 Alice 有一条消息 $m$ 要发送给 Bob，她可以对共享密钥进行 hash \(h = \mathcal{H}(S)\)，计算 $x = m + h$，并将 $x$ 发送给 Bob。Bob 计算 $h' = \mathcal{H}(S)$，计算 $m = x - h'$，从而获得 $m$。
\end{enumerate}   

外部观察者由于"Diffie-Hellman 问题"（DHP）将无法轻易计算共享密钥，DHP 表示从 $K_A$ 和 $K_B$ 找到 $S$ 非常困难。此外，DLP 阻止了他们找到 $k_A$ 或 $k_B$。\footnote{DHP 被认为至少与 DLP 具有相当的难度，尽管尚未被证明。\cite{diffie-hellman-problem}}


\subsection{Schnorr 签名和 Fiat-Shamir 变换 (Schnorr signatures and the Fiat-Shamir transform)}
\label{sec:schnorr-fiat-shamir}

1989年，Claus-Peter Schnorr 发表了一个现在著名的交互式认证协议 \cite{schnorr-signatures}，后由 Maurer 在2009年进行了推广 \cite{simple-zk-proof-maurer}，该协议允许某人证明自己知道给定公钥 $K$ 的私钥 $k$ 而不泄露任何关于它的信息 \cite{Signatures2015BorromeanRS}。它大致如下：

\begin{enumerate}
	\item 证明者生成一个随机整数 \(\alpha \in_R \mathbb{Z}_l\)，\footnote{\label{notation3_note}符号说明：\(\alpha \in_R \mathbb{Z}_l\) 中的 $R$ 表示 $\alpha$ 是从 \(\{1,2,3,...,l-1\}\) 中随机选取的。\marginnote{src/crypto/ crypto.cpp {\tt random32\_ unbiased()}} 换句话说，$\mathbb{Z}_l$ 是所有整数$\pmod l$。我们排除了 `$l$'，因为无穷远点在这里没有用。}计算 $\alpha G$，并将 $\alpha G$ 发送给验证者。
	\item 验证者生成一个随机{\em 挑战} $c \in_R \mathbb{Z}_l$ 并将 $c$ 发送给证明者。
	\item 证明者计算{\em 响应} $r = \alpha + c*k$ 并将 $r$ 发送给验证者。
	\item 验证者计算 $R = r G$ 和 $R' = \alpha G + c*K$，并检查 $R \stackrel{?}{=} R'$。
\end{enumerate}

验证者可以在证明者之前计算 $R' = \alpha G + c*K$，所以提供 $c$ 就像是在说，"我挑战你用 $R'$ 的离散对数来回应。"这是一个只有知道 $k$ 的证明者才能克服的挑战（除了可忽略的概率外）。

如果 $\alpha$ 是由证明者随机选择的，那么 $r$ 是随机分布的 \cite{SCOZZAFAVA1993313}，并且 $k$ 在 $r$ 中是信息论安全的（仍然可以通过对 $K$ 或 $\alpha G$ 求解 DLP 来找到它）。\footnote{\label{information_theoretic_note}具有信息论安全性的密码系统是指即使拥有无限计算能力的对手也无法破解它，因为他们根本没有足够的信息。}然而，如果证明者重复使用 $\alpha$ 来证明他对 $k$ 的知识，任何知道 $r = \alpha + c*k$ 和 $r' = \alpha + c'*k$ 中两个挑战的人都可以计算 $k$（两个方程，两个未知数）。\footnote{如果证明者是一台计算机，你可以想象有人在计算机生成 $\alpha$ 后"克隆"/复制该计算机，然后向每个副本提出不同的挑战。}\vspace{.175cm}%在安全证明中，诸如"分叉引理"和"成功时回退"等短语类似于这种克隆攻击。例如参见 \cite{Liu2004}。
\[k = \frac{r-r'}{c-c'}\]

如果证明者从一开始就就知道 $c$（例如，如果验证者秘密地给了她），她可以生成一个随机响应 $r$ 并计算 $\alpha G = r G - c K$。当她后来将 $r$ 发送给验证者时，她在从未需要知道 $k$ 的情况下"证明"了对 $k$ 的知识。观察证明者和验证者之间事件记录的人完全不会知情。该方案不是{\em 公开可验证的}。\cite{Signatures2015BorromeanRS}\\

在验证者作为挑战者的角色中，他在接收到 $\alpha G$ 后输出一个随机数，使他等价于一个{\em 随机函数}。随机函数，如 hash 函数，被称为随机预言机，因为计算一个就像是在向某人请求一个随机数 \cite{Signatures2015BorromeanRS}。\footnote{更一般地，``[在密码学中……预言机是任何可以提供关于系统的额外信息的系统，而这些信息否则无法获得。''\cite{cryptographic-oracle}}\\

使用 hash 函数（而非验证者）来生成挑战被称为 {\em Fiat-Shamir 变换} \cite{fiat-shamir-transform}，因为它使交互式证明变为非交互式且公开可验证 \cite{Signatures2015BorromeanRS}。\footnote{密码学 hash 函数 $\mathcal{H}$ 的输出在可能的输出范围内均匀分布。也就是说，对于某个输入 $A$，$\mathcal{H}(A) \in^D_R \mathbb{S}_H$，其中 $\mathbb{S}_H$ 是 $\mathcal{H}$ 的可能输出集合。我们使用 $\in^D_R$ 表示该函数是确定性随机的。$\mathcal{H}(A)$ 每次产生相同的结果，但其输出等价于一个随机数。
}\footnote{注意，非交互式类 Schnorr 证明（和签名）需要使用固定的生成元 $G$，或将生成元包含在挑战 hash 中。包含生成元的方式被称为密钥前缀，我们在后面会进一步讨论（第 \ref{blsag_note} 和 \ref{sec:robust-key-aggregation} 节）。}

\subsubsection*{非交互式证明 (Non-interactive proof)}

\begin{enumerate}
	\item 生成随机数 $\alpha \in_R \mathbb{Z}_l$，并计算 $\alpha G$。
	\item 使用密码学安全 hash 函数计算挑战，\(c = \mathcal{H}([\alpha G])\)。
	\item 定义响应 $r = \alpha + c*k$。
	\item 公开证明对 $(\alpha G, r)$。
\end{enumerate}

\subsubsection*{验证 (Verification)}

\begin{enumerate}
	\item 计算挑战：\(c' = \mathcal{H}([\alpha G])\)。
	\item 计算 $R = r G$ 和 $R' = \alpha G + c'*K$。
	\item 如果 $R = R'$，则证明者必须知道 $k$（除了可忽略的概率外）。
\end{enumerate}

\subsubsection*{为什么有效 (Why it works)}

\begin{align*}
r G &= (\alpha + c*k) G \\
	&= (\alpha G) + (c*k G) \\
	&= \alpha G + c*K \\
  R &= R'
\end{align*}

任何证明/签名方案的一个重要部分是验证它们所需的资源。这包括存储证明的空间和验证所花费的时间。在此方案中，我们存储一个 EC 点和一个整数，并需要知道公钥——另一个 EC 点。由于 hash 函数计算速度相对较快，请记住验证时间主要是椭圆曲线运算的函数\marginnote{src/ringct/ rctOps.cpp}。


\subsection{对消息签名 (Signing messages)}
\label{sec:signing-messages}

通常，密码学签名是对消息的密码学 hash 而非消息本身执行的，这有助于对不同大小的消息进行签名。然而，在本报告中，我们将宽泛地使用术语"消息"及其符号 $\mathfrak{m}$ 来指代消息本身和/或其 hash 值，除非另有说明。

对消息签名是互联网安全的基础，它让消息接收者确信其内容与签名者的意图一致。一种常见的签名方案叫做 ECDSA。参见 \cite{ecdsa}、ANSI X9.62 和 \cite{Hankerson:2003:GEC:940321}。

我们在此提出的签名方案是之前变换后的 Schnorr 证明的另一种表述。以这种方式思考签名为我们在下一章探索环签名做好了准备。

\subsubsection*{签名 (Signature)}

假设 Alice 拥有私钥/公钥对 \((k_A, K_A)\)。为了明确地签署任意消息 $\mathfrak{m}$，她可以执行以下步骤：

\begin{enumerate}
	\item 生成随机数 $\alpha \in_R \mathbb{Z}_l$，并计算 $\alpha G$。
	\item 使用密码学安全 hash 函数计算挑战，\(c = \mathcal{H}(\mathfrak{m},[\alpha G])\)。
	\item 定义响应 $r$ 使得 $\alpha = r + c*k_A$。换句话说，$r = \alpha - c*k_A$。
	\item 公开签名 $(c, r)$。
\end{enumerate}

\subsubsection*{验证 (Verification)}

任何知道 EC 域参数（指定使用了哪条椭圆曲线）、签名 $(c,r)$ 和签名方法、$\mathfrak{m}$ 和 hash 函数以及 $K_A$ 的第三方都可以验证签名：

\begin{enumerate}
	\item 计算挑战：\(c' = \mathcal{H}(\mathfrak{m},[r G + c*K_A])\)。
	\item 如果 $c = c'$，则签名通过。
\end{enumerate}

在此签名方案中，我们存储两个标量，并需要一个公钥。

\subsubsection*{为什么有效 (Why it works)}

这源于以下事实
\begin{align*}
 	r G &= (\alpha - c*k_A) G \\
 		&= \alpha G - c*K_A \\
\alpha G &= r G + c*K_A \\
\mathcal{H}_n(\mathfrak{m},[\alpha G]) &= \mathcal{H}_n(\mathfrak{m},[r G + c*K_A]) \\
       c &= c'
\end{align*}

因此 $k_A$ 的所有者（Alice）为 $\mathfrak{m}$ 创建了 $(c,r)$：她对消息进行了签名。其他人——一个没有 $k_A$ 的伪造者——能够制作 $(c,r)$ 的概率可以忽略不计，因此验证者可以确信消息没有被篡改。



%----------CURVE ED25519
\section{Ed25519 曲线 (Curve Ed25519)}
\label{Ed25519_section}

Monero 使用一种特定的扭曲 Edwards 椭圆曲线进行密码学运算，即 {\em Ed25519}，它是 Montgomery 曲线 {\em Curve25519} 的{\em 双有理等价}曲线\footnote{\label{birational_note}不作进一步详述，双有理等价可以被理解为可用有理项表达的同构。}。

Curve25519 和 Ed25519 均由 Bernstein {\em 等人} 发布 \cite{Bernstein2008, Bernstein2012, Bernstein2007}。\footnote{Bernstein 博士还开发了一种名为 ChaCha 的加密方案 \cite{Bernstein_chacha,chacha-irtf}，Monero 的主要实现使用它来加密\marginnote{src/wallet/ ringdb.cpp}与用户钱包相关的某些敏感信息。}

该\marginnote{src/crypto/ crypto\_ops\_ builder/ ref10Comm- entedComb- ined/\\\\ ge.h}曲线定义在素域 \(\mathbb{F}_{2^{255} - 19}\) 上（即 $q = 2^{255}-19$），通过以下方程：\vspace{.175cm}
\[-x^2 + y^2 = 1 - \frac{121665}{121666} x^2 y^2\]

这条曲线解决了密码学社区提出的许多问题。\footnote{即使一条曲线表面上没有密码学安全问题，创建它的人/组织也可能知道一个秘密问题，该问题只在非常罕见的曲线中出现。这样的人可能需要随机生成许多曲线，才能找到一条具有隐藏弱点且没有已知弱点的曲线。如果需要对曲线参数提供合理的解释，那么找到密码学社区可以接受的弱曲线就更加困难了。Ed25519 曲线被称为"完全刚性"曲线，这意味着其生成过程得到了充分解释。\cite{elliptic-curve-rigidity}}众所周知，NIST\footnote{\label{NIST_note}美国国家标准与技术研究院，\url{https://www.nist.gov/}}
标准算法存在问题。例如，最近已清楚 NIST 标准随机数生成算法 PNRG（基于椭圆曲线的版本）存在缺陷并包含潜在的后门 \cite{hales2014nsa}。从更广泛的视角来看，NIST 等标准化机构导致了密码学的单一文化，引入了中心化问题。一个很好的例子是 NSA 利用其对 NIST 的影响力来削弱国际密码学标准 \cite{NSA-NIST}。

Ed25519 曲线不受任何专利约束（参见 \cite{ECC-patents} 关于此主题的讨论），其背后的团队已经
开发\marginnote{src/crypto/ crypto\_ops\_ builder/}并调整了以效率为导向的基本密码学算法 \cite{Bernstein2007}。

扭曲 Edwards 曲线的阶可表示为 \(N=2^c l\)，其中 \(l\) 是素数，\(c\) 是正整数。在 Ed25519 曲线的情况下，其阶是一个76位数（$l$ 为253位）：\footnote{这意味着 Ed25519 中的私钥为253位。}\vspace{.175cm}
\[2^3 \cdot 7237005577332262213973186563042994240857116359379907606001950938285454250989\marginnote{src/ringct/ rctOps.h {\tt curve- Order()}}\]


\subsection{二进制表示 (Binary representation)}
\label{binary_note}
\(\mathbb{F}_{2^{255} - 19} \) 的元素被编码为256位整数，因此可以用32字节表示。由于每个元素只需要255位，最高有效位始终为零。

因此，Ed25519 中的任何点都可以用64字节表示。然而，通过应用下面描述的{\em 点压缩}技术，可以将此量减半至32字节。


\subsection{点压缩 (Point compression)}
\label{point_compression_section}

Ed25519 曲线具有其点可以被轻松压缩的特性，使得表示一个点仅消耗一个坐标的空间。我们将不深入证明这一特性所需的数学，但可以简要介绍其工作原理 \cite{eddsa-ed25519-irtf}。Ed25519 曲线的点压缩首次在 \cite{Bernstein2012} 中描述，而该概念首次在 \cite{Miller:point-compression-origin} 中引入。

这种点压缩方案源于扭曲 Edwards 曲线方程的变换（假设 $a = -1$，这对 Monero 成立）：$x^2 = (y^2-1)/(d y^2+1)$，\footnote{这里 $d = - \frac{121665}{121666}$。}这表明对于每个 $y$ 有两个可能的 $x$ 值（$+$ 或 $-$）。域元素 $x$ 和 $y$ 是$\pmod{q}$计算的，因此没有实际的负值。然而，对 $–x$ 取$\pmod{q}$ 会在奇偶之间改变值，因为 $q$ 是奇数。例如：$3 \pmod{5} = 3$，$-3 \pmod{5} = 2$。换句话说，域元素 $x$ 和 $–x$ 具有不同的奇偶分配。

如果我们有一个曲线点并知道其 $x$ 是偶数，但给定其 $y$ 值，变换后的曲线方程输出一个奇数，那么我们就知道对该数取反将给出正确的 $x$。一位可以传达这一信息，而方便的是 $y$ 坐标有一个额外的位。

假设我们想要压缩一个点 \((x, y)\)。

\begin{description}
	\item[编码 (Encoding)] 我们\marginnote{src/crypto/ crypto\_ops\_ builder/ ref10Comm- entedComb- ined/\\\\ ge\_to- bytes.c}将 $y$ 的最高有效位设为：如果 $x$ 是偶数则为0，如果是奇数则为1。得到的值 $y'$ 将代表该曲线点。
	
	\item[解码 (Decoding)] \hfill
	    \begin{enumerate}
    	    \item 获取\marginnote{ge\_from- bytes.c}[2.05cm]压缩后的点 $y'$，然后将其最高有效位复制到奇偶位 $b$，再将其设为0。现在它又变回了原始的 $y$。
    	    \item 设 \(u = y^2-1 \pmod q\) 和 \(v = d y^2  + 1 \pmod q\)。这意味着 $x^2 = u/v \pmod q$。
    	    \item 计算\footnote{由于 $q = 2^{255}-19 \equiv 5 \pmod{8}$，$(q-5)/8$ 和 $(q-1)/4$ 是整数。} \(z = u v^3 (u v^7)^{(q-5)/8} \pmod q\)。
            \begin{enumerate}
                \item 如果 \(v z^2 = u \pmod q\) 则 \(x' = z\)。
                \item 如果 \(v z^2 = -u \pmod q\) 则计算 \(x' = z*2^{(q-1)/4} \pmod q\)。
            \end{enumerate}
            \item 使用第一步中的奇偶位 \(b\)，如果 $b \ne$ $x'$ 的最低有效位，则 \(x = -x' \pmod q\)，否则 \(x = x'\)。
            \item 返回解压后的点 $(x,y)$。
	    \end{enumerate}
\end{description}

Ed25519 的实现（如 Monero）通常使用生成元 $G = (x,4/5)$ \cite{Bernstein2012}，其中 x 是基于 \(y = 4/5 \pmod q\) 点解压的"偶数"或 $b = 0$ 变体。


\subsection{EdDSA 签名算法 (EdDSA signature algorithm)}
\label{EdDSA_section}

Bernstein 和他的团队基于 Ed25519 曲线开发了许多基本算法。\footnote{\label{group_ops_twisted_edwards_note}参见\marginnote{src/crypto/ crypto\_ops\_ builder/ ref10Comm- entedComb- ined/} \cite{Bernstein2007} 了解扭曲 Edwards EC 中的高效群运算（即点加法、倍乘、混合加法等）。参见 \cite{curve25519} 了解高效的模运算。}
为了说明，我们将描述一种高度优化和安全的 ECDSA 签名方案替代方案，据作者称，使用普通的 Intel Xeon 处理器每秒可产生超过100,000个签名 \cite{Bernstein2012}。该算法也可以在互联网 RFC8032 \cite{rfc8032} 中找到描述。注意这是一种非常类似 Schnorr 的签名方案。

除了其他方面，它不使用每次生成随机整数，而是使用从签名者私钥和消息本身派生的 hash 值。这规避了与随机数生成器实现相关的安全缺陷。此外，该算法的另一个目标是避免访问秘密或不可预测的内存位置，以防止所谓的{\em 缓存时序攻击} \cite{Bernstein2012}。

我们在此提供该算法执行步骤的概述。完整的描述和 Python 语言的示例实现可在 \cite{rfc8032} 中找到。

\subsubsection*{签名 (Signature)}

\begin{enumerate}
	\item 设 \(h_k\) 为签名者私钥 \(k\) 的 hash \(\mathcal{H}(k)\)。
	计算 \(\alpha\) 为私钥 hash 和消息的 hash \(\alpha = \mathcal{H}(h_k, \mathfrak{m})\)。根据实现，$\mathfrak{m}$ 可以是实际消息或其 hash \cite{rfc8032}。
	
	\item 计算 \(\alpha G\) 和挑战 $ch = \mathcal{H}([\alpha G], K,  \mathfrak{m})$。

	\item 计算响应 \(r = \alpha + ch \cdot k\)。
	
	\item 签名为对 \((\alpha G, r)\)。
\end{enumerate}

\subsubsection*{验证 (Verification)}
验证按如下方式执行

\begin{enumerate}
	\item 计算 \(ch' = \mathcal{H}([\alpha G], K,  \mathfrak{m})\)。
	
	\item 如果等式 \(2^c r G \stackrel{?}{=} 2^c \alpha G + 2^c ch'*K \) 成立，则签名有效。
\end{enumerate}

$2^c$ 项来自 Bernstein {\em 等人} 的 EdDSA 算法通用形式 \cite{Bernstein2012}。根据该论文，虽然对于充分验证来说并非必需，但去除 $2^c$ 可以提供更强的方程。

公钥 $K$ 可以是任何 EC 点，但我们只想使用生成元 $G$ 的子群中的点。乘以余因子 $2^c$ 可确保所有点都在该子群中。或者，验证者可以检查 $l K \stackrel{?}{=} 0$，这仅在 $K$ 在子群中时才有效。我们不知道这些预防措施能捕获的任何弱点，不过正如我们将看到的，后一种方法在 Monero 中很重要（第 \ref{blsag_note} 节）。

在此签名方案中，我们存储一个 EC 点和一个标量，并有一个公钥。

\subsubsection*{为什么有效 (Why it works)}
\begin{align*}
2^c r G &= 2^c (\alpha + \mathcal{H}([\alpha G], K,  \mathfrak{m}) \cdot k) \cdot G \\
		&= 2^c \alpha G + 2^c \mathcal{H}([\alpha G], K,  \mathfrak{m}) \cdot K 
\end{align*}

\subsubsection*{二进制表示 (Binary representation)}

默认情况下，EdDSA 签名需要 \(64 + 32\) 字节用于 EC 点 $\alpha G$ 和标量 $r$。然而，RFC8032 假设点 \(\alpha G\) 被压缩，这将空间需求减少到仅 \(32 + 32\) 字节。我们包含公钥 $K$，这意味着总共 \(32 + 32 + 32\) 字节。



\section{二元运算符 XOR (Binary operator XOR)}
\label{sec:XOR_section}

二元运算符 XOR 是一个有用的工具，将在第 \ref{sec:integrated-addresses} 和 \ref{sec:pedersen_monero} 节中出现。它接受两个参数，如果其中一个（但不是两个）为真则返回真 \cite{wolfram-xor}。以下是其真值表：

\begin{center}
    \begin{tabular}{|c|c|c|}
    \hline
        A & B & A XOR B \\
    \hline\hline
        T & T & F \\
    \hline
        T & F & T \\
    \hline
        F & T & T \\
    \hline
        F & F & F \\
    \hline
    \end{tabular}
\end{center}

在计算机科学的语境中，XOR 等价于模2的位加法。例如，两对位的 XOR：
\begin{alignat*}{1}
    \text{XOR}(\{1,1\},\{1,0\}) &= \{1+1,1+0\} \pmod 2 \\
                                &= \{0,1\} 
\end{alignat*}

以下也产生 $\{0,1\}$：$\text{XOR}(\{1,0\},\{1,1\})$、$\text{XOR}(\{0,0\},\{0,1\})$ 和 $\text{XOR}(\{0,1\},\{0,0\})$。对于具有 $b$ 位的 XOR 输入，有 $2^{\text{b}} - 1$ 种其他输入组合可以产生相同的输出。这意味着如果 $C = \text{XOR}(A,B)$ 且输入 $A \in_R \{0,...,2^{\text{b}-1}\}$，知道 $C$ 的观察者不会获得关于 $B$ 的任何信息。

同时，任何知道 $\{A,B,C\}$ 中两个元素的人（其中 $C = \text{XOR}(A,B)$），都可以计算第三个元素，例如 $A = \text{XOR}(B,C)$。XOR 指示两个元素是不同还是相同，所以知道 $C$ 和 $B$ 就足以暴露 $A$。仔细检查真值表可以揭示这一关键特性。\footnote{XOR 的一个有趣应用（与 Monero 无关）是在不使用第三个寄存器的情况下交换两个位寄存器。我们使用符号 $\oplus$ 表示 XOR 运算。$A \oplus A = 0$，所以在寄存器之间进行三次 XOR 运算后：$\{A, B\} \rightarrow{} \{[A \oplus B], B\} \rightarrow{} \{[A \oplus B], B \oplus [A \oplus B]\} = \{[A \oplus B], A \oplus 0\} = \{[A \oplus B], A\} \rightarrow{} \{[A \oplus B] \oplus A, A\} = \{B, A\}$。}
